% ============================================================
%  ch12  第五讲 §5.1 Satake 范畴、§5.2 $\Sat_G$ 作为 Tannaka 范畴
%  原文行 2464–2643
% ============================================================

\section{几何 Satake 等价}\label{Lecture V}

本讲讨论几何 Satake 等价，它是几何朗兰兹纲领的基石之一。除另有声明外，本节假设 $k$ 代数闭。设 $\mO=k[[t]]$、$F=k((t))$，$G$ 是 $k$ 上的约化群，$\underline G= G\otimes_k\mO$。若群 $G$ 明确，为简洁把 $\Gr_G$ 记作 $\Gr$。

固定一个与 $k$ 的特征不同的素数 $\ell$。层指 $\Ql$-层，所有推前与拉回都是\emph{导出的}。需要指出，几何 Satake 等价还可在另外两种层范畴中表述：(1) 若 $\cha\ k=0$，可用 D-模（见 \cite[\S\ 5.3]{BD}）；(2) 若 $k=\bC$，可用 $\Gr$ 上解析拓扑下 $E$-向量空间层，其中 $E$ 是特征零的域。下面给出的途径对这两种范畴同样适用。另一方面，几何 Satake 等价还有 $\bZ_\ell$-与 $\bF_\ell$-层版本，或者当 $k=\bC$ 时 $\Lambda$-模层版本，其中 $\Lambda$ 是有限整体维数的交换 noether 环（例如 $\Lambda=\bZ$）。此时下述途径不够用，一般系数的推广需参见 \cite{MV}。

还要指出，$p$-进群也有几何 Satake 等价，细节见 \cite{Z14}。

\subsection{Satake 范畴 $\Sat_G$}\label{Satake category}

首先，$L^+G$ 在 $\Gr$ 上的作用满足 \S\ref{A:ind} 中的条件 (S)：每个轨道形如 $\Gr_\mu=L^+G/(L^+G\cap t^{\mu} L^+G t^{-\mu})$，且 $L^+G\cap t^{\mu} L^+G t^{-\mu}$ 连通。因此 \S\ref{A:ind} 的构造给出 $\Gr$ 上 $L^+G$-等变反常层范畴 $\on{P}_{L^+G}(\Gr)$。以 $\IC_\mu$ 记 $\Gr_{\leq \mu}$ 上的相交上同调层。则 $\IC_\mu|_{\Gr_\mu}=\Ql[(2\rho,\mu)]$，且它在每个层片 $\Gr_\la$ 上的限制是常值的。又因 $L^+G\cap t^{\mu} L^+G t^{-\mu}$ 连通，$\on{P}_{L^+G}(\Gr)$ 的不可约对象恰为这些 $\on{IC}_\mu$。\zy{反常层之名有误导：nLab 词条 perverse sheaf（\url{https://ncatlab.org/nlab/show/perverse+sheaf}，2026-10-16 修订）开篇即指出这类对象「既非 perverse 也非 sheaf」，而 $\on{P}(X)$ 是 $D^b_c(X)$ 的全 Abel 子范畴。$\IC_\mu$ 与相交上同调的关系见 nLab 词条 intersection cohomology（2023-05-07 修订）。[来源：nLab perverse sheaf、intersection cohomology。]}

\begin{prop}\label{semisimplicity}
范畴 $\on{P}_{L^+G}(\Gr)$ 半单。
\end{prop}

\begin{proof}
要证对任意 $\la,\mu$ 有 $\on{Ext}^1(\IC_\la,\IC_\mu)=0$。本质情形是 $\la\leq \mu$ 或 $\mu\leq \la$。此时消没由以下两点推出：(i)（仿射）Schubert 簇的 $\on{IC}$-层茎上同调的纯性（参见 \cite{KL}）；(ii) $\Gr$ 每个连通分支中 $L^+G$-轨道维数的奇偶性（见引理 \ref{parity}）。细节见 \cite[Proposition 1]{Ga}（另见 \cite[Proposition 3.1]{Ri}）。
\end{proof}

由此得到

\begin{cor}
$\on{P}_{L^+G}(\Gr)$ 中的对象都是 $\Aut^+(D)$-等变的。
\end{cor}

\begin{proof}
因为每个 IC 层都必须是 $\Aut^+(D)$-等变的（由引理 \ref{A:forgetful}）。
\end{proof}

$\on{P}_{L^+G}(\Gr)$ 上有自然的幺半结构，由 Lusztig 的层卷积积给出。即对 $\mA_1,\mA_2\in \on{P}_{L^+G}(\Gr)$，以 $\mA_1\tilde\boxtimes\mA_2$ 记 $\mA_1$ 与 $\mA_2$ 在 $\Gr\tilde\times\Gr$ 上的「外扭曲积」（见 \S\ref{A:twist product} 与 \S\ref{A:ind}）。例如若 $\mA_i=\IC_{\mu_i}$，则 $\IC_{\mu_1}\tilde\boxtimes\IC_{\mu_2}$ 典范同构于 $\Gr_{\leq \mu_1}\tilde\times\Gr_{\leq\mu_2}$ 的相交上同调层。于是 $\mA_1$ 与 $\mA_2$ 的卷积积定义为
\begin{equation}\label{Lus conv}
\mA_1\star\mA_2:=m_!(\mA_1\tilde\boxtimes\mA_2),
\end{equation}
其中 $m:\Gr\tilde\times\Gr\to \Gr$ 是卷积映射（由 \eqref{conv m0} 定义）。这是 $\Gr$ 上的 $L^+G$-等变 $\ell$-进复形。类似地，用 $n$ 重卷积映射 \eqref{conv m} 可定义 $n$ 重卷积积 $\mA_1\star\cdots\star\mA_n=m_!(\mA_1\tilde\boxtimes\cdots\tilde\boxtimes\mA_n)$。注意因 $m$ 归纳 proper，故 $m_!=m_*$。

\begin{prop}\label{conv prod}
卷积积 $\mA_1\star\mA_2$ 反常。
\end{prop}

\cite{BD} 称之为奇迹定理。它使我们能停留在 Abel 幺半范畴中，而不必处理更麻烦的三角幺半范畴概念。$\Sat_G$ 的另一个奇迹是引理 \ref{left=right}。

该命题本质上以两种不同途径证明。第一种基于仿射 Grassmann 簇的某种半无限几何（下面 \S\ref{MV semiinfinite} 将引入），Lusztig 用它证明了仿射 Hecke 代数的一个数值结果 \cite{Lu}，该结果与此命题等价（见 \cite[Proposition 2.2.1]{Gi1}）；Mirkovi\'c--Vilonen 用它证明了卷积映射的半小性，这与命题等价（见 \cite[Lemma 4.3, 4.4]{MV} 与下面注 \ref{semismall}）；\cite[Corollary 9.7]{NP}（另见 \cite[\S\ 2.2]{Z14}）也用它给出该命题的直接证明。

第二种途径基于 Beilinson--Drinfeld Grassmann 簇。Gaitsgory 借助邻近圈函子证明了更强的结果：在命题 \ref{conv prod} 中 $\mA_1$ 可以是 $\Gr$ 上任意反常层（注意定义外扭曲积并不需要 $\mA_1$ 的等变性条件）。见下面注 \ref{conv perverse}。

需要命题 \ref{conv prod} 的一个推论。设 $\mmu=(\mu_1,\ldots,\mu_n)$ 是支配余权序列。由命题 \ref{semisimplicity} 与 \ref{conv prod}，可写
\[\IC_{\mu_1}\star\cdots\star\IC_{\mu_n}=\bigoplus_{\la} V_{\mmu}^\la\otimes \IC_\la,\]
其中 $V_{\mmu}^\la=\Hom_{\on{P}(\Gr)}(\IC_\la, \IC_{\mu_1}\star\cdots\star\IC_{\mu_n})$。令 $|\mmu|=\sum \mu_i$。

\begin{cor}\label{Satake fiber}
有典范同构
\[V_{\mmu}^\la\cong \on{H}^{(2\rho,|\mmu|-\la)}_c(\Gr_{\mmu}\cap m^{-1}(t^\la),\Ql),\]
且后者有由 $\Gr_{\mmu}\cap m^{-1}(t^\la)$ 的维数为 $(2\rho,|\mmu|-\la)$ 的不可约分支给出的基。
\end{cor}

\begin{proof}
令 $d=\dim \Gr_{\mmu}\cap m^{-1}(t^\la)$。把 $\Gr_{\leq \mu_\bullet}$ 分层为 $\sqcup_{\mu'_\bullet}\Gr_{\mu'_\bullet}$，一个谱序列论证表明 $\IC_{\mu_1}\star\cdots\star\IC_{\mu_n}$ 在 $t^\la$ 处次数 $2d- (2\rho,|\mmu|)$ 的茎上同调为
\[\on{H}^{2d}_c(\Gr_{\mmu}\cap m^{-1}(t^\la),\Ql).\]
由 $\IC_{\mu_1}\star\cdots\star\IC_{\mu_n}$ 的反常性，$2d- (2\rho,|\mmu|)\leq -(2\rho,\la)$，等价于 $d\leq (\rho,|\mmu|-\la)$。又 $V_{\mmu}^\la$ 等于 $\IC_{\mmu}$ 在 $t^\la$ 处次数 $-(2\rho,\la)$ 的茎上同调。推论得证。
\end{proof}

\begin{rmk}\label{semismall}
(i) 注意该论证还表明
\[\dim  \Gr_{\leq \mmu}\cap m^{-1}(t^{\la})\leq (\rho,|\mmu|-\la).\]
于是卷积映射半小。但一般来说，$\Gr_{\leq \mmu}\cap m^{-1}(t^{\la})$ 可能存在维数严格小于 $(\rho,|\mmu|-\la)$ 的不可约分支，例子见 \cite[Remark 4.3, 8.3]{Ha}。

(ii) 这些 $\Gr_{\leq \mmu}\cap m^{-1}(t^{\la})$ 有时称为 \textbf{Satake 纤维}。当 $G=\GL_n$ 时，它们有时可等同于 Springer 纤维（见例 \ref{conv v.s. springer}）。特别地，这蕴含 Springer 消解 $\tilde\mN_n\to \mN_n$ 的半小性。
\end{rmk}

把 $(\mA_1\star\mA_2)\star\mA_3$ 与 $\mA_1\star(\mA_2\star\mA_3)$ 都等同于 $\mA_1\star\mA_2\star\mA_3$，就赋予卷积积自然结合约束。称幺半范畴 $(\on{P}_{L^+G}(\Gr),\star)$ 为 \textbf{Satake 范畴}，为简洁有时记作 $\Sat_G$。

\subsection{$\Sat_G$ 作为 Tannaka 范畴}\label{mon str of H}

本小节先按 \cite[\S\ 2.3]{Z14} 赋予超上同调函子 $\on{H}^*: \Sat_G\to \on{Vect}_{\Ql}$ 幺半结构，再说明如何把 $(\Sat_G,\on{H}^*)$ 上的幺半结构升级为对称幺半结构。真正的证明将基于后面对卷积积作为融合积的解释。\zy{这里用的是 Tannaka 对偶的「纤维函子」形式：由张量函子 $F:\mC\to\on{Vect}$ 的张量自同构群重建群概形，一般形式见 nLab 词条 Tannaka duality（\url{https://ncatlab.org/nlab/show/Tannaka+duality}，2025-07-28 修订）；本文把它用于 $\mC=\Sat_G$。[来源：nLab Tannaka duality，2025-07-28。]}

\begin{prop}\label{fiber functor}
超上同调函子 $\on{H}^*:\Sat_G\to \on{Vect}_{\Ql}$ 有典范幺半结构。
\end{prop}

想法如下。先设 $k=\bC$，从而可如 \S\ref{top loop group} 那样把 $\Gr$ 与 $\Gr\tilde\times\Gr$ 看作无穷维复解析空间。注意范畴 $\on{P}_{L^+G}(\Gr_G)$ 也可用解析拓扑定义。容易验证下述事实（见 \cite[Proposition 2.1.1]{Gi1}）。

\begin{lem}
在定理 \ref{top triv} 的同构下，$\Gr\tilde\times\Gr$ 上的 $\mA_1\tilde\boxtimes\mA_2$ 等同于 $\Omega K\times\Omega K$ 上的 $\mA_1\boxtimes\mA_2$。
\end{lem}

由 K\"unneth 公式可知 $\on{H}^*(-):=\on{H}^*(\Gr,-):\Sat_G\to \on{Vect}_{\Ql}$ 具有命题 \ref{fiber functor} 所述的幺半结构。

若 $k$ 一般，定理 \ref{top triv} 与上述引理并不字面成立。有两种途径使这些拓扑同构有意义。一种基于 Beilinson--Drinfeld Grassmann 簇把 $\Gr\tilde\times \Gr$ 形变到 $\Gr\times \Gr$，将在 \S\ \ref{fusion} 解释。这里解释另一种用等变上同调的途径，其优点是完全局部（即不依赖曲线的选取）。

回忆对 $\mA\in\Sat_G$，其 $L^+G$-等变上同调 $\on{H}^*_{L^+G}(\mA)$ 是自由 $R_{G}$-模（参见 \S\ref{A:char class}）。把 $\Gr$ 上 $L^+G$-等变层（非正式地）看作 $L^+G\backslash LG/L^+G$ 上的层，则 $\on{H}^*_{L^+G}(\mA)$ 上应还有另一个 $R_{G}$-模结构；它可以如下字面构造。以 $L^+G^{(m)}\subset L^+G$ 记第 $m$ 个同余子群（即投影 $L^+G\to L^mG$ 的核），以 $\Gr^{(m)}=LG/L^+G^{(m)}$ 记 $\Gr$ 上的通用 $L^m G$-主丛。则 $\Gr^{(m)}$ 带有 $L^+G\times L^mG$ 的作用，且投影
\[\pi_m:\Gr^{(m)}\to \Gr\]
是 $L^+G$-等变的。于是由 \eqref{coh free action}，
\[\on{H}^*_{L^+G}(\mA)= \on{H}^*_{L^+G\times L^mG}(\pi_m^*\mA),\]
从而给 $\on{H}^*_{L^+G}(\mA)$ 一个 $R_{G}$-双模结构。注意只要 $m>0$，该结构与 $m$ 无关。

回忆（含幺）代数 $R$ 的双模范畴有标准幺半结构 $(M,N)\mapsto M\otimes_RN$。

\begin{lem}\label{soergel monoidal}
对任意 $\mA_1,\mA_2,\ldots,\mA_n$，存在 $R_{G}$-双模的典范同构
\[\on{H}^*_{L^+G}(\mA_1\star\cdots\star\mA_n)\cong \on{H}^*_{L^+G}(\mA_1)\otimes_{R_{G}}\cdots\otimes_{R_{G}}\on{H}^*_{L^+G}(\mA_n),\]
并满足自然的相容条件。因此 $\on{H}^*_{L^+G}(-):\Sat_G\to (\on{R}_{G}\otimes\on{R}_{G})\on{-mod}$ 上有自然幺半结构。
\end{lem}

这是 Soergel 想法的一个变体：把（仿射）旗簇上的等变层对应到 Soergel 双模。这里只解释 $n=2$、$\mA_i=\IC_{\mu_i}$ 时如何构造映射，一般情形类似（例如见 \cite[\S~2.3]{Z14}）。设 $L^+G$ 在 $\Gr_{\leq \mu_2}$ 上的作用通过 $L^mG$ 分解，于是 $\Gr_{\leq \mu_1}\tilde\times\Gr_{\leq \mu_2}= \Gr_{\leq\mu_1}^{(m)}\times^{L^mG}\Gr_{\leq \mu_2}$。由卷积积的定义与 \eqref{coh free action}，
\[\on{H}^*_{L^+G}(\IC_{\mu_1}\star\IC_{\mu_2})=\on{H}^*_{L^+G\times L^mG}(\pi_m^*\IC_{\mu_1}\boxtimes \IC_{\mu_2}).\]
另一方面，由等变 K\"unneth 公式 \eqref{equiv Kunneth}，有典范同构
\[\on{H}^*_{L^+G}(\IC_{\mu_1})\otimes_{R_{G}} \on{H}^*_{L^+G}(\IC_{\mu_2})\cong \on{H}^*_{L^+G\times L^mG}(\pi_m^*\IC_{\mu_1}\boxtimes \IC_{\mu_2}).\]
合起来便得所需同构。

还需要一个事实。

\begin{lem}\label{left=right}
$\on{H}^*_{L^+G}(\mA)$ 上的两个 $R_{G}$-结构相同。
\end{lem}

\begin{proof}
这里只说明思路，细节见 \cite[\S\ 2.3]{Z14}。用引理 \ref{minusch} 与引理 \ref{quasiminusch}，容易对 $\mA=\IC_\mu$（$\mu$ 极小或拟极小）证明本引理。再由引理 \ref{soergel monoidal}，结论对这些 $\IC_\mu$ 的卷积成立。最后，每个 $\mA\in\Sat_G$ 都作为这些卷积的直和项出现（这是所谓 PRV 猜想的几何版本，见 \cite[Proposition 9.6]{NP}）。引理得证。
\end{proof}

\begin{rmk}
本引理是 $\underline G=G\otimes \mO$ 情形所特有的。例如它对仿射旗簇 $\Fl_G$ 上的等变反常层不成立（术语见注 \ref{terminology}）。
\end{rmk}

注意把命题 \ref{soergel monoidal} 与 \ref{left=right} 结合，实际上赋予等变超上同调函子
\[\on{H}^*_{L^+G}: \Sat_G\to \on{Proj}_{R_G}\]
典范幺半结构，其中如前 $\on{Proj}_{R_G}$ 记有限投射 $R_G$-模范畴，配以通常张量积。由 \eqref{de-equiv}，沿增广映射 $R_{G}\to \Ql$ 特殊化后，便得 $\on{H}^*$ 上所需的幺半结构。

留给读者验证：当 $k=\bC$ 时，上述 $\on{H}^*$ 上的幺半结构与用同胚 $\Gr\tilde\times\Gr\cong \Gr\times \Gr$ 构造者一致。

赋予 $\on{H}^*$ 幺半结构后，下面说明 $\Sat_G$ 上交换约束的构造，它使 $\on{H}^*$ 成为对称幺半函子。关键陈述如下。

\begin{prop}\label{existence of comm}
对任意 $\mA_1,\mA_2\in\Sat_G$，存在唯一同构 $c'_{\mA_1,\mA_2}:\mA_1\star\mA_2\cong\mA_2\star\mA_1$ 使下图交换
\[\begin{CD}
\on{H}^*(\mA_1\star\mA_2)@>\on{H}^*(c'_{\mA_1,\mA_2})>>\on{H}^*(\mA_2\star\mA_1)\\
@V\cong VV@VV\cong V\\
\on{H}^*(\mA_1)\otimes\on{H}^*(\mA_2)@>\cong>c_{\on{gr}}>\on{H}^*(\mA_2)\otimes\on{H}^*(\mA_1),
\end{CD}\]
其中竖直同构来自命题 \ref{monoidal functor}，而 $c_{\on{gr}}$ 是分次向量空间的交换约束，即
\[c_{\on{gr}}(v\otimes w)=(-1)^{\deg(v)\deg(w)}(w\otimes v).\]
\end{prop}

\begin{rmk}
读者应注意：命题 \ref{fiber functor} 给 $(\Sat_G, \on{H}^*)$ 添了额外结构，而命题 \ref{existence of comm} 建立的是 $(\Sat_G,\on{H}^*)$ 的某些性质。
\end{rmk}

因 $\on{H}^*:\Sat_G\to\on{Vect}_{\Ql}$ 忠实（由 $\Sat_G$ 半单），$c'_{\mA_1,\mA_2}$ 的唯一性显然，内容在其存在性。它将在把卷积积解释为融合积之后于 \S\ref{fusion} 证明。注意 \cite[\S~2.4]{Z14} 提示本命题可能存在纯局部（即不用整体曲线）的途径。

下面是所需的推论。首先可按 Koszul 符号改变来修改约束 $c'_{\mA_1,\mA_2}$（另见 \cite[\S~5.3.21]{BD}；更优雅的处理见 \cite{MV} 的注 6.2 之后）。即：范畴 $\on{P}_{L^+G}(\Gr)$ 带有由 \eqref{parity map} 诱导的 $\bZ/2$-分次。称 $\mA$ 具有纯奇偶性，若 $p(\on{Supp}(\mA))$ 为 $1$ 或 $-1$，此时定义 $p(\mA)=p(\on{Supp}(\mA))$，其中 $\on{Supp}(\mA)$ 记 $\mA$ 的支集。于是当 $\mA_1,\mA_2$ 具有纯奇偶性 $p(\mA_1),p(\mA_2)$ 时，令
\begin{equation}\label{koszul sign change}
c_{\mA_1,\mA_2}:=(-1)^{p(\mA_1)p(\mA_2)}c'_{\mA_1,\mA_2}.
\end{equation}
用修改后的交换约束，下图交换
\[\begin{CD}
\on{H}^*(\mA_1\star\mA_2)@>\on{H}^*(c_{\mA_1,\mA_2})>>\on{H}^*(\mA_2\star\mA_1)\\
@V\cong VV@VV\cong V\\
\on{H}^*(\mA_1)\otimes\on{H}^*(\mA_2)@>\cong>c_{\on{vect}}>\on{H}^*(\mA_2)\otimes\on{H}^*(\mA_1),
\end{CD}\]
其中同构 $c_{\on{vect}}$ 是向量空间的通常交换约束，即 $c_{\on{vect}}(v\otimes w)=w\otimes v$。

\begin{thm}\label{Tannakian structure}
幺半范畴 $\Sat_G$ 连同上述约束 $c_{\mA_1,\mA_2}$ 构成对称幺半范畴。超上同调函子 $\on{H}^*$ 是张量函子。事实上，$(\on{P}_{L^+G}(\Gr),\star,c,\on{H}^*)$ 构成中性 Tannaka 范畴。
\end{thm}

\begin{proof}
需验证 $c_{\mA_2,\mA_1}c_{\mA_1,\mA_2}=\id$ 与六边形公理。由 $\on{H}^*$ 的忠实性，取上同调后证明即可。由命题 \ref{existence of comm} 与 $c_{\on{vect}}^2=\id$ 知 $\on{H}^*(c_{\mA_2,\mA_1}c_{\mA_1,\mA_2})=\id$，故 $c_{\mA_2,\mA_1}c_{\mA_1,\mA_2}=\id$。六边形公理类似可证。第二个陈述显然。

为证 $(\on{P}_{L^+G}(\Gr),\star,c,\on{H}^*)$ 是中性 Tannaka 范畴，应用 \cite[Proposition 1.20, Remark 2.18]{DM}。即：仍由半单性，$\on{H}^*$ 正合（且忠实）。容易验证\emph{同上}命题 1.20 的条件 (a)--(d)（单位对象取 $\IC_0$）。至于条件 (e)，注意 $\dim \on{H}^*(\Gr,\IC_\mu)=1$ 当且仅当 $\dim \Gr_{\leq \mu}=0$，此时 $\mu$ 是中心 $Z_G\subset G$ 的余特征；于是 $\IC_{-\mu}\star\IC_\mu\cong \IC_0$。
\end{proof}
